\documentclass[conference]{IEEEtran}
\IEEEoverridecommandlockouts

\usepackage{cite}
\usepackage{amsmath,amssymb,amsfonts}
\usepackage{graphicx}
\usepackage{textcomp}
\usepackage{xcolor}
\usepackage{hyperref}

% ICSE 2027 NIER — limite typique : 4 pages (corps), références non comprises.
% Vérifier le call for papers officiel avant soumission finale pour la limite exacte.

\begin{document}

\title{Native-IA Software Engineering: From Empirical Synthesis to a Validated Governance Framework\\
\large{[Titre provisoire — à affiner Phase 1]}}

\author{
  \IEEEauthorblockN{A. Yvan GUIFO FODJO}
  \IEEEauthorblockA{
    EFREI Paris / Sorbonne Université (LIP6) \\
    Paris, France \\
    [email à compléter]
  }
}

\maketitle

\begin{abstract}
[Rédiger après stabilisation de l'introduction — Phase 1, S1-S2.
L'abstract NIER doit tenir en ~150-200 mots et annoncer clairement :
(1) le problème, (2) le gap, (3) la contribution, (4) la validation envisagée.]
\end{abstract}

\begin{IEEEkeywords}
AI-native software engineering, governance, traceability, semantic validation,
human-AI collaboration, empirical software engineering
\end{IEEEkeywords}

\section{Introduction}
\label{sec:introduction}

% Semaine 1-2 (Phase 1) : positionner le NIER par rapport au papier #184.
% Le NIER doit soit (a) être un extrait resserré et reformulé du framework Native-IA
% validé empiriquement par les projets de la formation AI Portfolio, soit (b) être une
% extension avec validation, selon l'issue de la review de #184.
%
% Points à couvrir :
%   - Contexte : montée des systèmes agentiques en génie logiciel
%   - Gap : absence de cadre de gouvernance opérationnel et validé empiriquement
%   - Contribution : framework + validation sur cas d'étude (Sprint 2-3 formation)
%   - Structure de l'article

[À rédiger — Phase 1]

\section{Related Work}
\label{sec:related-work}

% Semaine 3 (Phase 1) : réutiliser et actualiser les références vérifiées du papier #184
% (voir references.bib). Organiser par thème :
%   - Qualité du code généré et dette technique
%   - Limites des outils de validation traditionnels
%   - Traçabilité, gouvernance, responsabilité
%   - Facteurs humains et cognitifs
% Ajouter les preprints arXiv pertinents détectés par la veille (scripts/veille/arxiv_monitor.py).

[À rédiger — Phase 1]

\section{Methodology}
\label{sec:methodology}

% Semaine 5 (Phase 2) : décrire la méthode de synthèse empirique (héritée de #184,
% inspirée PRISMA) ET le protocole de validation sur cas d'étude (Enterprise AI
% Assistant, Sprint 3 formation — Phase 3).

[À rédiger — Phase 2]

\section{Framework}
\label{sec:framework}

% Semaine 6 (Phase 2) : le coeur du papier. Trois dimensions (héritées de #184) :
%   1. Gouvernance explicite des interactions humain-IA
%   2. AI Commit Ledger (provenance des contributions)
%   3. Validation sémantique intégrée au pipeline de développement
% NOUVEAU pour le NIER : intégrer les observations empiriques de LangGraph
% (multi-agents), LangSmith (observabilité/traçabilité), sécurité LLM (validation).

[À rédiger — Phase 2]

\section{Discussion}
\label{sec:discussion}

% Semaine 7 (Phase 2) : trade-offs velocity/quality/control, autonomie/contrôle,
% délégation cognitive/maintien des compétences. Threats to validity (hérité #184,
% à actualiser avec le cas d'étude).

[À rédiger — Phase 2]

\section{Conclusion and Future Work}
\label{sec:conclusion}

% Phase 4 (semaines 13-15) : conclusion finale + perspectives.
% Mentionner explicitement : standardisation traçabilité/validation, étude empirique
% étendue des impacts cognitifs, conception de processus pour environnements
% hautement automatisés.

[À rédiger — Phase 4]


\bibliographystyle{IEEEtran}
\bibliography{references}

\end{document}
